\documentclass[JEP,XML,SOM,Unicode,published]{cedram}
\datereceived{2022-11-03}
\dateaccepted{2023-03-21}
\dateepreuves{2023-03-27}

\usepackage[scr=rsfs,cal=euler]{mathalfa}
\usepackage[all,cmtip]{xy}\let\labelstyle\textstyle
\newcommand{\Changel}{\mathcode`l="7160}
\newcommand{\Changelback}{\mathcode`l="716C}
\newcommand{\NoChangel}[1]{%
\expandafter\let\csname old\string#1\endcsname=#1
\let#1=\relax
\newcommand{#1}{\mathcode`l="716C\csname old\string#1\endcsname\mathcode`l="7160 }%
}

\NoChangel{\log}
\NoChangel{\ln}
\NoChangel{\lim}
\NoChangel{\limsup}
\NoChangel{\liminf}
\NoChangel{\varlimsup}
\NoChangel{\varliminf}
\NoChangel{\varinjlim}
\NoChangel{\varprojlim}

\let\oldbigsqcup\bigsqcup
\renewcommand{\bigsqcup}{\mathop{\textstyle\oldbigsqcup}\displaylimits}
\newcommand{\spfrac}[2]{\sfrac{#1}{(#2)}}
\newcommand\mto{\mathchoice{\longmapsto}{\mapsto}{\mapsto}{\mapsto}}

\theoremstyle{plain}
\newtheorem*{maintheorem}{Main Theorem}
\newtheorem{theorem}{Theorem}
\newtheorem{lemma}[theorem]{Lemma}
\newtheorem{proposition}[theorem]{Proposition}

\theoremstyle{definition}
\newtheorem{definition}[theorem]{Definition}

\theoremstyle{remark}
\newtheorem{remark}[theorem]{Remark}

\DeclareMathOperator{\Gal}{Gal}\NoChangel{\Gal}

\datepublished{2023-03-31}
\begin{document}
\frontmatter
\title{Rational maps with rational multipliers}

\author[\initial{V.} \lastname{Huguin}]{\firstname{Valentin} \lastname{Huguin}}
\address{Constructor University Bremen gGmbH\\
Campus Ring 1, 28759 Bremen, Germany}
\email{vhuguin@constructor.university}
\urladdr{}
\thanks{The research of the author was supported by the German Research Foundation (DFG, project number 455038303)}

\subjclass{37P35, 37F10, 37P05}
\keywords{Multiplier spectrum, exceptional maps, arithmetic dynamics}

\begin{abstract}
In this article, we show that every rational map whose multipliers all lie in a given number field is a power map, a Chebyshev map or a Lattès map. This strengthens a conjecture by Milnor concerning rational maps with integer multipliers, which was recently proved by Ji and Xie.
\end{abstract}

\altkeywords{Multiplicateur, application exceptionnelle, dynamique arithmétique}

\alttitle{Fractions rationnelles avec multiplicateurs rationnels}

\begin{altabstract}
Dans cet article, nous montrons que toute fraction rationnelle dont les multiplicateurs sont tous dans un corps de nombres donné est une application puissance, une application de Tchebychev ou un exemple de Lattès. Ceci généralise une conjecture de Milnor concernant les fractions rationnelles avec multiplicateurs entiers, qui a été récemment démontrée par Ji et Xie.
\end{altabstract}
\maketitle
\tableofcontents
\mainmatter
\Changel

\section{Introduction}

Suppose that $f \colon \widehat{\mathbb{C}} \to \widehat{\mathbb{C}}$ is a rational map. A point $z_{0} \in \widehat{\mathbb{C}}$ is said to be \emph{periodic} for~$f$ if there exists an integer $p \geq 1$ such that $f^{\circ p}(z_{0}) = z_{0}$. In this case, the least such integer $p$ is called the \emph{period} of $z_{0}$ and $\left\lbrace f^{\circ n}(z_{0}) : n \geq 0 \right\rbrace$ is said to be a \emph{cycle} for $f$. The \emph{multiplier} of $f$ at $z_{0}$ is the unique eigenvalue $\lambda \in \mathbb{C}$ of the differential of~$f^{\circ p}$ at~$z_{0}$, so that $\lambda = \left(f^{\circ p} \right)^{\prime}(z_{0})$ if $z_{0} \in \mathbb{C}$. The map $f$ has the same multiplier at each point of the cycle. Furthermore, the multiplier is invariant under conjugation: if~$\phi \colon \widehat{\mathbb{C}} \to \widehat{\mathbb{C}}$ is a Möbius transformation and $g = \phi \circ f \circ \phi^{-1}$, then $\phi(z_{0})$ is periodic for $g$ with period $p$ and multiplier $\lambda$.

The notion of multiplier is fundamental in complex dynamics and plays a major role in the study of both local and global dynamics of rational maps. Furthermore, multipliers almost determine rational maps up to conjugation: aside from flexible Lattès maps, there are only finitely many conjugacy classes of rational maps that have the same collection of multipliers at each period (see~\cite[Cor.\,2.3]{MM1987}).

The purpose of this article is to show that the multipliers of a rational map do not all lie in a given number field unless the rational map is exceptional.

\begin{definition}
A rational map $f \colon \widehat{\mathbb{C}} \to \widehat{\mathbb{C}}$ of degree $d \geq 2$ is said to be a \emph{power map} if it is conjugate to $z \mto z^{\pm d}$.
\end{definition}

For every $d \geq 2$, there exists a unique polynomial $T_{d} \in \mathbb{C}[z]$ such that
\[
T_{d}\left(z +z^{-1} \right) = z^{d} +z^{-d}.
\]
The polynomial $T_{d}$ is monic of degree $d$ and is called the $d$th \emph{Chebyshev polynomial}.

\begin{definition}
A rational map $f \colon \widehat{\mathbb{C}} \to \widehat{\mathbb{C}}$ of degree $d \geq 2$ is said to be a \emph{Chebyshev map} if it is conjugate to $\pm T_{d}$.
\end{definition}

Power maps and Chebyshev maps are exactly the finite quotients of affine maps on cylinders (see~\cite[Lem.\,3.8]{M2006}). There are also finite quotients of affine maps on tori.

\begin{definition}
A rational map $f \colon \widehat{\mathbb{C}} \to \widehat{\mathbb{C}}$ of degree $d \geq 2$ is said to be a \emph{Lattès map} if there exist a torus $\mathbb{T} = \mathbb{C}/\Lambda$, with $\Lambda$ a lattice in $\mathbb{C}$, a holomorphic map $L \colon \mathbb{T} \to \mathbb{T}$ and a nonconstant holomorphic map $p \colon \mathbb{T} \to \widehat{\mathbb{C}}$ that make the following diagram commute:
\[
\xymatrix{
\mathbb{T}\ar[r]^-L\ar[d]_p&\mathbb{T}\ar[d]^p\\
\widehat{\mathbb{C}}\ar[r]^-f&\widehat{\mathbb{C}}
}
\]
Furthermore, if $p$ has degree $2$ and $L$ is of the form
\[
L \colon z +\Lambda \mto a z +b +\Lambda, \quad \text{with} \quad a \in \mathbb{Z}, \, b \in \mathbb{C},
\]
then $f$ is said to be a \emph{flexible} Lattès map.
\end{definition}

Power maps, Chebyshev maps and Lattès maps play a special role in complex dynamics and are sometimes called exceptional. For example, their multipliers all lie in a discrete subring of $\mathbb{C}$.

An \emph{imaginary quadratic field} is a number field of the form $\mathbb{Q}(i \sqrt{D})$, with $D$ a positive integer. Given a number field $K$, we denote by $\mathcal{O}_{K}$ its ring of integers.

\skpt
\begin{proposition}[{\cite[Cor.\,3.9 \& Lem.\,5.6]{M2006}}]
\label{proposition:special}
\begin{itemize}
\item Suppose that $f \colon \widehat{\mathbb{C}} \to \widehat{\mathbb{C}}$ is a power map or a Chebyshev map. Then $f$ has only integer multipliers.
\item Suppose that $f \colon \widehat{\mathbb{C}} \to \widehat{\mathbb{C}}$ is a Lattès map. Then there exists an imaginary quadratic field $K$ such that the multipliers of $f$ all lie in $\mathcal{O}_{K}$. Furthermore, the multipliers of $f$ are all integers if and only if $f$ is flexible.
\end{itemize}
\end{proposition}

Here, we are interested in the converse of Proposition~\ref{proposition:special}. We prove the following:

\begin{maintheorem}
Assume that $K$ is a number field and $f \colon \widehat{\mathbb{C}} \to \widehat{\mathbb{C}}$ is a rational map of degree $d \geq 2$ whose multipliers all lie in $K$. Then $f$ is a power map, a Chebyshev map or a Lattès map.
\end{maintheorem}

Main Theorem strengthens several recent results. These were motivated by a conjecture made by Milnor in~\cite{M2006}, which states that power maps, Chebyshev maps and flexible Lattès maps are the only rational maps that have only integer multipliers. This problem was first studied by the author in~\cite{H2021}, where the conjecture was proved for unicritical polynomial maps and cubic polynomial maps with symmetries. More generally, in light of Proposition~\ref{proposition:special}, one could ask whether power maps, Chebyshev maps and Lattès maps are the only rational maps whose multipliers all lie in the ring of integers of a given imaginary quadratic field. In~\cite{H2022}, the author gave a positive answer to this question for quadratic rational maps. Finally, the general case was settled by Ji and Xie, who thus proved Milnor's conjecture.

\begin{theorem}[{\cite[Th.\,1.13]{JX2022}}]
\label{theorem:milnor}
Assume that $K$ is an imaginary quadratic field and $f \colon \widehat{\mathbb{C}} \to \widehat{\mathbb{C}}$ is a rational map of degree $d \geq 2$ whose multipliers all lie in $\mathcal{O}_{K}$. Then $f$ is a power map, a Chebyshev map or a Lattès map.
\end{theorem}

In fact, the author proved in~\cite{H2021} that every unicritical polynomial map that has only rational multipliers is either a power map or a Chebyshev map. Thus, one could even ask whether Theorem~\ref{theorem:milnor} still holds if $K$ is an arbitrary number field and the multipliers of $f$ are only assumed to lie in $K$. Main Theorem provides a positive answer to this question.

We note that Ji and Xie's proof of Theorem~\ref{theorem:milnor} makes crucial use of the fact that the multipliers all lie in a discrete subring of $\mathbb{C}$, whereas the set of multipliers may a priori have limit points under the assumption of Main Theorem. Therefore, we take a different approach. Our proof of Main Theorem relies on an equidistribution result for points of small height proved by Autissier in~\cite{A2001} and on a characterization of power maps, Chebyshev maps and Lattès maps proved by Zdunik in~\cite{Z2014}.

Finally, let us mention that Eremenko and van~Strien investigated the rational maps that have only real multipliers in~\cite{EvS2011}. They proved that a rational map $f \colon \widehat{\mathbb{C}} \to \widehat{\mathbb{C}}$ of degree $d \geq 2$ has this property if and only if it is a flexible Lattès map or its Julia set $\mathcal{J}_{f}$ is contained in a circle.

\subsubsection*{Acknowledgements}
The author would like to thank Xavier Buff, Thomas Gauthier, Igors Gorbovickis and Jasmin Raissy for interesting discussions about this work and the referees for their suggestions.

\section{Proof of the result}

First, let us state the equidistribution result and the characterization of power maps, Chebyshev maps and Lattès maps that we use in our proof of Main Theorem.

We denote by $\overline{\mathbb{Q}}$ the field of algebraic numbers. If $f \colon \widehat{\mathbb{C}} \to \widehat{\mathbb{C}}$ is a rational map of degree $d \geq 2$ defined over $\overline{\mathbb{Q}}$, then its periodic points all lie in $\overline{\mathbb{Q}} \cup \lbrace \infty \rbrace$. Given a number field $K$, we extend each $\sigma \in \Gal\left(\overline{\mathbb{Q}}/K \right)$ to a map $\sigma \colon \overline{\mathbb{Q}} \cup \lbrace \infty \rbrace \to \overline{\mathbb{Q}} \cup \lbrace \infty \rbrace$ by setting $\sigma(\infty) = \infty$.

Given a rational map $f \colon \widehat{\mathbb{C}} \to \widehat{\mathbb{C}}$ of degree $d \geq 2$, we denote by $\mu_{f}$ its measure of maximal entropy, which is a Borel probability measure on $\widehat{\mathbb{C}}$ whose support is the Julia set $\mathcal{J}_{f}$ of $f$. We refer the reader to~\cite{FLM1983}, \cite{L1983} and~\cite{M1983} for further information.

Recall that a sequence $\left(\mu_{n} \right)_{n \geq 0}$ of Borel probability measures on $\widehat{\mathbb{C}}$ \emph{converges weakly} to a Borel probability measure $\nu$ on $\widehat{\mathbb{C}}$ if
\[
\lim_{n \to +\infty} \int_{\widehat{\mathbb{C}}} \varphi \, d\mu_{n} = \int_{\widehat{\mathbb{C}}} \varphi \, d\nu
\]
for all continuous functions $\varphi \colon \widehat{\mathbb{C}} \to \mathbb{R}$.

The statement below is a particular case of an equidistribution result for points of small height due to Autissier. We refer to~\cite[Th.\,10.24]{BR2010} and~\cite[Th.\,2 \& Th.\,4]{FRL2006} for an analogous result in a dynamical context that explicitly implies the statement below.

\begin{theorem}[{\cite[Prop.\,4.1.4]{A2001}}]
\label{theorem:equidistribution}
Assume that $f \colon \widehat{\mathbb{C}} \to \widehat{\mathbb{C}}$ is a rational map of degree $d \geq 2$ defined over a number field $K$ and $\left(S_{n} \right)_{n \geq 0}$ is a sequence of nonempty, pairwise disjoint and $\Gal\left(\overline{\mathbb{Q}}/K \right)$\nobreakdash-invariant sets of periodic points for $f$. For $n \geq 0$, define the Borel probability measure
\[
\mu_{n} = \frac{1}{\left\lvert S_{n} \right\rvert} \sum_{z \in S_{n}} \delta_{z}.
\]
Then $\left(\mu_{n} \right)_{n \geq 0}$ converges weakly to $\mu_{f}$.
\end{theorem}

Given a rational map $f \colon \widehat{\mathbb{C}} \to \widehat{\mathbb{C}}$ and a periodic point $z_{0} \in \widehat{\mathbb{C}}$ for $f$ with period $p$ and multiplier $\lambda$, the \emph{characteristic exponent} of $f$ at $z_{0}$ is
\[
\chi_{f}(z_{0}) = \frac{1}{p} \log\lvert \lambda \rvert.
\]

Suppose that $f \colon \widehat{\mathbb{C}} \to \widehat{\mathbb{C}}$ is a rational map of degree $d \geq 2$. We denote by $\left\lVert f^{\prime} \right\rVert$ the norm of the differential of $f$ with respect to the spherical metric $\spfrac{2 \lvert dz \rvert}{1 +\lvert z \rvert^{2}}$; this is the unique continuous function $\left\lVert f^{\prime} \right\rVert \colon \widehat{\mathbb{C}} \to \mathbb{R}_{\geq 0}$ that satisfies
\[
\left\lVert f^{\prime}(z) \right\rVert = \frac{\left\lvert f^{\prime}(z) \right\rvert \left(1 +\lvert z \rvert^{2} \right)}{1 +\left\lvert f(z) \right\rvert^{2}}
\]
for all $z \in \mathbb{C}$ such that $f(z) \in \mathbb{C}$. The \emph{Lyapunov exponent} of $f$ is
\[
\mathcal{L}_{f} = \int_{\widehat{\mathbb{C}}} \log\left\lVert f^{\prime} \right\rVert \, d\mu_{f},
\]
where $\mu_{f}$ denotes the measure of maximal entropy of $f$.

In our proof of Main Theorem, we use the following characterization of power maps, Chebyshev maps and Lattès maps:

\begin{theorem}[{\cite[Prop.\,4]{Z2014}}]
\label{theorem:exceptional}
Assume that $f \colon \widehat{\mathbb{C}} \to \widehat{\mathbb{C}}$ is a rational map of degree $d \geq 2$ that is not a power map, a Chebyshev map or a Lattès map. Then $f$ has infinitely many periodic points with characteristic exponent greater than $\mathcal{L}_{f}$.
\end{theorem}

\begin{remark}
In fact, Zdunik only states that every rational map $f \colon \widehat{\mathbb{C}} \to \widehat{\mathbb{C}}$ of degree $d \geq 2$ that is not a power map, a Chebyshev map or a Lattès map has at least one periodic point with characteristic exponent greater than $\mathcal{L}_{f}$, but her proof can be easily modified to obtain the statement above. As Theorem~\ref{theorem:exceptional} plays a key role in our proof of Main Theorem, let us explain how Zdunik's proof can be changed. Fix an integer $N \geq 1$. On~\cite[p.\,260]{Z2014}, denote by $p_{1}, \dotsc, p_{s}$ not only the critical values for $f^{M}$ but also the periodic points for $f$ with period at most $N$ in order to obtain a family $\mathcal{B}$ of balls that also contain no periodic point for $f$ with period less than or equal to $N$. Then the remainder of~\cite[Proof of Prop.\,4]{Z2014} shows that there exist a ball $B \in \mathcal{B}$ and a periodic point $z_{0} \in B$ for $f$ such that $\chi_{f}(z_{0}) > \mathcal{L}_{f}$. The point $z_{0}$ has period greater than $N$ by the construction of $\mathcal{B}$. Since this holds for every $N \geq 1$, the map $f$ has infinitely many periodic points with characteristic exponent greater than $\mathcal{L}_{f}$.
\end{remark}

Given a rational map $f \colon \widehat{\mathbb{C}} \to \widehat{\mathbb{C}}$, the \emph{postcritical set} of $f$ is
\[
\mathcal{P}_{f} = \bigcup_{n \geq 1} f^{\circ n}(\mathcal{C}_{f}),
\]
where $\mathcal{C}_{f}$ denotes the set of critical points for $f$.

We also need the following approximation lemma:

\begin{lemma}
\label{lemma:homoclinic}
Suppose that $f \colon \widehat{\mathbb{C}} \to \widehat{\mathbb{C}}$ is a rational map of degree $d \geq 2$ and $z_{0} \in \widehat{\mathbb{C}}$ is a repelling periodic point for $f$ that does not lie in $\mathcal{P}_{f}$. Then there is a sequence $(w_{n})_{n \geq 0}$ of periodic points for $f$ with pairwise distinct periods such that
\[
\lim_{n \to +\infty} \chi_{f}(w_{n}) = \chi_{f}(z_{0}).
\]
\end{lemma}

\begin{proof}
Replacing $f$ by an iterate if necessary, we may assume that $z_{0}$ is a repelling fixed point for $f$. Then there exist connected open neighborhoods $U$ and $V$ of $z_{0}$ such that $\overline{U} \subset V$ and $f$ has no critical point in $\overline{U}$ and induces a biholomorphism $f_{U} \colon U \to V$. Moreover, $z_{0}$ lies in the Julia set $\mathcal{J}_{f}$ of $f$, its iterated preimages form a dense subset of $\mathcal{J}_{f}$ and $\mathcal{J}_{f}$ has no isolated point, and hence there exist $l \geq 1$ and $z_{l} \in V \setminus \left\lbrace z_{0} \right\rbrace$ such that $f^{\circ l}\left(z_{l} \right) = z_{0}$. The point $z_{l}$ is not critical for $f^{\circ l}$ because $z_{0}$ does not lie in $\mathcal{P}_{f}$ by hypothesis, and hence there is an open neighborhood $W$ of $z_{l}$ such that $\overline{W} \subset V$ and $f^{\circ l}$ has no critical point in $\overline{W}$ and induces a biholomorphism $\left(f^{\circ l} \right)_{W} \colon W \to f^{\circ l}(W)$. Denote by $g \colon V \to U$ and $h \colon f^{\circ l}(W) \to W$ the inverses of~$f_{U}$ and $\left(f^{\circ l} \right)_{W}$, respectively. The map $g$ induces a contraction of $\overline{U}$ with respect to the Poincaré metric on $V$ as $\overline{U} \subset V$. Therefore, $z_{0}$ is the unique periodic point for~$g$ and there exists an integer $N \geq l$ such that $g^{\circ (n -l)}(V) \subset f^{\circ l}(W)$ for all $n \geq N$ since $g(z_{0}) = z_{0}$ and $f^{\circ l}(W)$ is a neighborhood of $z_{0}$. For $n \geq N$, we can consider $h \circ g^{\circ (n -l)} \colon V \to W$. For every $n \geq N$, the map $h \circ g^{\circ (n -l)}$ induces a contraction of $\overline{W}$ with respect to the Poincaré metric on $V$ since $\overline{W} \subset V$, and hence it has a unique fixed point $w_{n} \in W$. Let us prove that $w_{n}$ is periodic for $f$ with period $n$ for all $n \geq \max\lbrace N, 2 l \rbrace$. For every $n \geq N$ and every $j \in \lbrace l, \dotsc, n \rbrace$, we have
\[
f^{\circ j}(w_{n}) = f^{\circ j} \circ h \circ g^{\circ (n -l)}(w_{n}) = g^{\circ (n -j)}(w_{n}).
\]
It follows that, for each $n \geq N$, the point $w_{n}$ is periodic for $f$ with period dividing $n$. Suppose now that $n \geq \max\lbrace N, 2 l \rbrace$ and the period of $w_{n}$ is a proper divisor $p_{n}$ of $n$. Then $n -p_{n} \in \lbrace l, \dotsc, n -1 \rbrace$ and $f^{\circ \left(n -p_{n} \right)}(w_{n}) = w_{n}$, and hence $w_{n}$ is periodic for $g$ by the relation above. Therefore, $w_{n} = z_{0}$ since $z_{0}$ is the unique periodic point for $g$, and hence
\[
z_{0} = h \circ g^{\circ (n -l)}(z_{0}) = h(z_{0}) = z_{l}
\]
by the definition of $w_{n}$, which is a contradiction. Thus, $w_{n}$ is periodic for $f$ with period $n$ for all $n \geq \max\lbrace N, 2 l \rbrace$. Finally, it remains to prove that
\[
\lim_{n \to +\infty} \chi_{f}(w_{n}) = \chi_{f}(z_{0}).
\]
For $n \geq \max\lbrace N, 2 l \rbrace$, denote by $\lambda_{n}$ the multiplier of $f$ at $w_{n}$, so that
\[
\left\lvert \lambda_{n} \right\rvert = \bigl\lVert(f^{\circ l})^{\prime}(w_{n}) \bigr\rVert \cdot \prod_{j = l}^{n -1} \left\lVert f^{\prime}\bigl(f^{\circ j}(w_{n}) \bigr) \right\rVert = \bigl\lVert(f^{\circ l})^{\prime}(w_{n}) \bigr\rVert \cdot \prod_{j = 1}^{n -l} \bigl\lVert f^{\prime}\bigl(g^{\circ j}(w_{n}) \bigr) \bigr\rVert.
\]
Suppose that $\alpha \in (1, +\infty)$. Since $\left\lVert f^{\prime} \right\rVert$ is continuous at $z_{0}$, there is a neighborhood~$O_{\alpha}$ of $z_{0}$ such that
\[
\frac{\exp\left(\chi_{f}(z_{0}) \right)}{\alpha} \leq \left\lVert f^{\prime}(z) \right\rVert \leq \alpha \exp\left(\chi_{f}(z_{0}) \right)
\]
for all $z \in O_{\alpha}$. As $g(z_{0}) = z_{0}$ and $g$ induces a contraction of $\overline{U}$ with respect to the Poincaré metric on $V$, there exists $J_{\alpha} \geq 1$ such that $g^{\circ j}(V) \subset O_{\alpha}$ for all $j \geq J_{\alpha}$. Define
\[
m = \min\Bigl\lbrace \min_{\overline{U}} \left\lVert f^{\prime} \right\rVert, \min_{\overline{W}} \bigl\lVert(f^{\circ l})^{\prime} \bigr\rVert \Bigr\rbrace \in \mathbb{R}_{> 0} \quad \text{and} \quad M = \max_{\widehat{\mathbb{C}}} \left\lVert f^{\prime} \right\rVert \in \mathbb{R}_{> 0}.
\]
For every $n \geq \max\left\lbrace N, 2 l, J_{\alpha} +l -1 \right\rbrace$, we have
\[
m^{J_{\alpha}} \Bigl(\frac{\exp\left(\chi_{f}(z_{0}) \right)}{\alpha} \Bigr)^{n -J_{\alpha} -l +1} \leq \left\lvert \lambda_{n} \right\rvert \leq M^{J_{\alpha} +l -1} \bigl(\alpha \exp(\chi_{f}(z_{0})) \bigr)^{n -J_{\alpha} -l +1},
\]
and hence
\[
\chi_{f}(w_{n}) \geq \frac{J_{\alpha} \log(m)}{n} +\frac{\left(n -J_{\alpha} -l +1 \right) \left(\chi_{f}(z_{0}) -\log(\alpha) \right)}{n}
\]
and
\[
\chi_{f}(w_{n}) \leq \frac{\left(J_{\alpha} +l -1 \right) \log(M)}{n} +\frac{\left(n -J_{\alpha} -l +1 \right) \left(\chi_{f}(z_{0}) +\log(\alpha) \right)}{n}.
\]
Therefore, we have
\[
\chi_{f}(z_{0}) -\log(\alpha) \leq \liminf_{n \to +\infty} \chi_{f}(w_{n}) \leq \limsup_{n \to +\infty} \chi_{f}(w_{n}) \leq \chi_{f}(z_{0}) +\log(\alpha).
\]
As this holds for all $\alpha \in (1, +\infty)$, the lemma is proved.
\end{proof}

\begin{remark}
On the one hand, the statement of Lemma~\ref{lemma:homoclinic} is rather weak. Using the Koenigs linearization theorem and the same technique as in the proof above, Ji and Xie show in~\cite[\S2]{JX2022} that, if $f \colon \widehat{\mathbb{C}} \to \widehat{\mathbb{C}}$ is a rational map of degree $d \geq 2$ and $z_{0} \in \widehat{\mathbb{C}} \setminus \mathcal{P}_{f}$ is a repelling fixed point for $f$ with multiplier $\lambda$, then there is a sequence $(w_{n})_{n \geq N}$ such that $w_{n}$ is periodic for $f$ with period $n$ and multiplier $\rho_{n}$ for all $n \geq N$ and $\rho_{n} = a \lambda^{n} +b +o(1)$ as $n \to +\infty$ for some $a \in \mathbb{C}^{*}$ and $b \in \mathbb{C}$.

On the other hand, our proof of Lemma~\ref{lemma:homoclinic} can be generalized by using several periodic points. Thus, one can show that, if $f \colon \widehat{\mathbb{C}} \to \widehat{\mathbb{C}}$ is a rational map of degree $d \geq 2$, then the closure of the set of characteristic exponents of $f$ at its cycles is the union of an interval and a finite set (compare~\cite[Proof of Lem.\,2.1]{FG2022}). This is a weak version of~\cite[Cor.\,1.16]{JX2022}, which also states that this interval is not a singleton if the map is not a power map, a Chebyshev map or a Lattès map.
\end{remark}

We now prove our result.

\begin{proof}[Proof of Main Theorem]
Assume that $f \colon \widehat{\mathbb{C}} \to \widehat{\mathbb{C}}$ is a rational map of degree $d \geq 2$ whose multipliers all lie in a given number field $K$. By~\cite[Th.\,2.1]{S1998}, the moduli space $\mathcal{M}_{d}(\mathbb{C})$ of rational maps of degree $d$ is an algebraic variety defined over $\mathbb{Q}$. Denote by $\mathcal{Z}_{f}$ the set of conjugacy classes $[g] \in \mathcal{M}_{d}(\mathbb{C})$ such that $f$ and $g$ have the same multipliers with the same multiplicities at their cycles with period $n$ for each $n \geq 1$. Then $\mathcal{Z}_{f}$ is a Zariski closed subset of $\mathcal{M}_{d}(\mathbb{C})$ defined over $K$ by~\cite[Th.\,4.5]{S1998}. Moreover, it consists of finitely many elements of $\mathcal{M}_{d}(\mathbb{C})$ and possibly a curve of conjugacy classes of flexible Lattès maps by~\cite[Cor.\,2.3]{MM1987}. Therefore, each rational map $g \colon \widehat{\mathbb{C}} \to \widehat{\mathbb{C}}$ of degree $d$ such that $[g] \in \mathcal{Z}_{f}$ is conjugate to a rational map defined over a finite extension of $K$ or is a flexible Lattès map, and this holds in particular for $f$. If $f$ is a flexible Lattès map, we are done. Thus, conjugating $f$ and replacing $K$ by a finite extension if necessary, we may assume that $f$ is defined over~$K$.

Suppose that $z_{0} \in \widehat{\mathbb{C}}$ is a repelling periodic point for $f$ that does not lie in $\mathcal{P}_{f}$, and let us prove that $\chi_{f}(z_{0}) \leq \mathcal{L}_{f}$. By Lemma~\ref{lemma:homoclinic}, there is a sequence $(w_{n})_{n \geq 0}$ of periodic points for $f$ with pairwise distinct periods such that
\[
\lim_{n \to +\infty} \chi_{f}(w_{n}) = \chi_{f}(z_{0}).
\]
For $n \geq 0$, denote by $p_{n}$ and $\lambda_{n}$ the period and multiplier of $w_{n}$ for $f$, respectively, and define
\[
S_{n} = \bigl\lbrace \sigma\left(f^{\circ j}(w_{n}) \right) : j \in\lbrace 0, \dotsc, p_{n} -1\rbrace, \, \sigma \in \Gal(\overline{\mathbb{Q}}/K) \bigr\rbrace
\]
to be the smallest $\Gal(\overline{\mathbb{Q}}/K)$\nobreakdash-invariant subset of $\overline{\mathbb{Q}} \cup \lbrace \infty \rbrace$ that contains the cycle for~$f$ containing $w_{n}$. As $f$ is defined over $K$,
\[
S_{n} = \bigcup_{\sigma \in \Gal(\overline{\mathbb{Q}}/K)} \left\lbrace f^{\circ j}\left(\sigma(w_{n}) \right) : j \in\lbrace 0, \dotsc, p_{n} -1\rbrace \right\rbrace
\]
is a union of cycles for $f$ with period $p_{n}$ for all $n \geq 0$. In particular, $S_{m} \cap S_{n} = \varnothing$ for all distinct $m, n \geq 0$ since $p_{m} \neq p_{n}$. For every $n \geq 0$ and every $\sigma \in \Gal(\overline{\mathbb{Q}}/K)$, denoting by $\rho_{\sigma}^{(n)}$ the multiplier of $f$ at $\sigma(w_{n})$, we have $\rho_{\sigma}^{(n)} = \sigma\left(\lambda_{n} \right)$ as $f$ is defined over $K$ and $\sigma(w_{n})$ has period $p_{n}$, and hence $\rho_{\sigma}^{(n)} = \lambda_{n}$ since $\lambda_{n} \in K$ by hypothesis. Therefore, for every $n \geq 0$, choosing $\sigma_{1}, \dotsc, \sigma_{r} \in \Gal(\overline{\mathbb{Q}}/K)$---which depend on $n$---such that
\[
S_{n} = \bigsqcup_{l = 1}^{r} \bigl\lbrace f^{\circ j}\left(\sigma_{l}(w_{n}) \right) : j \in \left\lbrace 0, \dotsc, p_{n} -1 \right\rbrace \bigr\rbrace,
\]
we have
\[
\frac{1}{\left\lvert S_{n} \right\rvert} \sum_{z \in S_{n}} \log\left\lVert f^{\prime}(z) \right\rVert = \frac{1}{\left\lvert S_{n} \right\rvert} \sum_{l = 1}^{r} \log\vert \rho_{\sigma_{l}}^{(n)}\vert = \frac{r}{\left\lvert S_{n} \right\rvert} \log\left\lvert \lambda_{n} \right\rvert = \chi_{f}(w_{n}).
\]
For $n \geq 0$, define the Borel probability measure
\[
\mu_{n} = \frac{1}{\left\lvert S_{n} \right\rvert} \sum_{z \in S_{n}} \delta_{z}.
\]
For every $m \in \mathbb{R}$ and every $n \geq 0$, we have
\[
\chi_{f}(w_{n}) \leq \frac{1}{\left\lvert S_{n} \right\rvert} \sum_{z \in S_{n}} \max\left\lbrace \log\lVert f^{\prime}(z)\rVert, m \right\rbrace = \int_{\widehat{\mathbb{C}}} \max\left(\log\lVert f^{\prime}\rVert, m \right) \, d\mu_{n}.
\]
By Theorem~\ref{theorem:equidistribution}, the sequence $\left(\mu_{n} \right)_{n \geq 0}$ converges weakly to $\mu_{f}$. Therefore, for every $m \in \mathbb{R}$, we obtain
\[
\chi_{f}(z_{0}) \leq \int_{\widehat{\mathbb{C}}} \max\left(\log\left\lVert f^{\prime} \right\rVert, m \right) \, d\mu_{f}
\]
by letting $n \to +\infty$ since $\max\left(\log\left\lVert f^{\prime} \right\rVert, m \right)$ is continuous on $\widehat{\mathbb{C}}$. Taking the limit as $m \to -\infty$, it follows from the monotone convergence theorem that $\chi_{f}(z_{0}) \leq \mathcal{L}_{f}$. Thus, we have shown that $\chi_{f}(z_{0}) \leq \mathcal{L}_{f}$ for all but finitely many periodic points $z_{0} \in \widehat{\mathbb{C}}$ for~$f$. Therefore, $f$ is a power map, a Chebyshev map or a Lattès map by Theorem~\ref{theorem:exceptional}, which completes the proof of the theorem.
\end{proof}

\Changelback

\backmatter
\bibliographystyle{jepalpha+eid}
\bibliography{huguin}
\end{document}